\chapter*{Resumen}
\addcontentsline{toc}{chapter}{Resumen}

Este informe presenta el diseño, la implementación y la evaluación de una plataforma web para gestionar el catálogo y los pedidos de un restaurante mediante persistencia políglota. MongoDB funciona como fuente operacional de verdad y Apache Cassandra mantiene dos proyecciones de actividad orientadas a consultas. La integración entre ambos almacenes utiliza una salida transaccional (\emph{transactional outbox}) y un trabajador idempotente con entrega al menos una vez. La arquitectura se compone de una interfaz React, una API NestJS, MongoDB y Cassandra, ejecutados en un entorno local de cuatro contenedores con recursos limitados.

El trabajo sigue una metodología guiada por especificaciones. Los requisitos se expresan mediante escenarios verificables, se relacionan con decisiones de diseño y se comprueban mediante pruebas unitarias, de integración y de extremo a extremo. El documento presta especial atención a la atomicidad del registro de pedidos, el cálculo de totales en el servidor, la validez de las transiciones de estado, la idempotencia, la recuperación ante fallos y el diseño de Cassandra a partir de patrones de consulta. La evaluación se apoya en artefactos y criterios de aceptación reproducibles, evitando depender únicamente de capturas de pantalla.

\noindent\textbf{Palabras clave:} persistencia políglota, MongoDB, Apache Cassandra, bases de datos NoSQL, salida transaccional, idempotencia, Docker Compose.

\chapter*{Abstract}
\addcontentsline{toc}{chapter}{Abstract}

This report presents the design, implementation, and evaluation of a restaurant catalog and order management platform based on polyglot persistence. MongoDB acts as the operational source of truth, while Apache Cassandra maintains two query-oriented activity projections. Both stores are integrated through a transactional outbox and an idempotent worker providing at-least-once delivery. The architecture consists of a React interface, a NestJS API, MongoDB, and Cassandra, running in a resource-bounded four-container local environment.

The project follows a specification-driven methodology. Requirements are stated as verifiable scenarios, linked to design decisions, and verified through unit, integration, and end-to-end testing. The report focuses on atomic order creation, server-side total calculation, valid state transitions, idempotency, failure recovery, and query-first Cassandra modeling. The evaluation uses reproducible artifacts and acceptance criteria instead of relying only on screenshots.

\noindent\textbf{Keywords:} polyglot persistence, MongoDB, Apache Cassandra, NoSQL databases, transactional outbox, idempotency, Docker Compose.
